\documentclass[11pt,a4paper]{article}
\usepackage[utf8]{inputenc}
\usepackage{amsmath,amssymb,amsfonts,amsthm}
\usepackage{graphicx}
\usepackage{booktabs}
\usepackage{algorithm}
\usepackage{algorithmic}
\usepackage{hyperref}
\usepackage{geometry}
\usepackage{microtype}
\usepackage{cite}

\geometry{margin=1in}

\title{\textbf{CrisisRoute: Uncertainty-Aware Humanitarian Logistics and Equitable Multi-Commodity Routing under Dynamic Road Passability}}

\author{
  \textbf{Farhan Aaqil}\\
  Department of Computer Science \& Engineering\\
  \texttt{fadurrani543@gmail.com}
}

\date{\today}

\newtheorem{theorem}{Theorem}
\newtheorem{lemma}{Lemma}
\newtheorem{definition}{Definition}

\begin{document}

\maketitle

\begin{abstract}
Post-disaster humanitarian logistics operations frequently suffer catastrophic delivery failures due to inaccurate road passability estimates, unquantified hazard uncertainty, and deterministic routing models that prioritize aggregate throughput over equitable survival. We present \textbf{CrisisRoute}, an end-to-end mathematical and operational decision-support system designed for rapid-onset disaster response. CrisisRoute integrates: (1) an isotonic-calibrated gradient-boosted passability estimator operating over dynamic 12-dimensional topological and meteorological features; (2) a risk-aware, non-linear penalized shortest-path routing algorithm ($c_e = t_e - \lambda_r \ln p_e + \lambda_s a_e$) generating diverse alternative corridors via Yen's $k$-shortest paths; (3) a two-stage stochastic recourse allocation model solved via Sample Average Approximation (SAA) and Constraint Programming (CP-SAT), enforcing non-negotiable service floors for vulnerable populations; and (4) an Expected Value of Information (VoI) engine directing field reconnaissance to links where uncertainty reduction maximizes operational regret minimization. We evaluate CrisisRoute across \textit{CrisisBench-mini}, a 240-instance benchmark comprising three flood and landslide corridors in India (Kerala Wayanad/Periyar, Assam Barak Valley, and Odisha Coastal). Empirical evaluations confirm that: post-hoc isotonic calibration reduces Expected Calibration Error (ECE) by 96.2\% ($p < 0.001$); risk-aware routing reduces route severance by 50.0\% ($p = 0.0012$) under ground-truth flash-flood cutoffs; CP-SAT achieves 100\% compliance with vulnerable population service floors while eliminating deprivation inequality (Gini $\le 0.15$); and VoI-guided drone reconnaissance yields a 25.0\% greater regret reduction than uniform inspection ($p < 0.001$). CrisisRoute establishes a rigorous, reproducible paradigm for uncertainty-quantified humanitarian operations.
\end{abstract}

\section{Introduction}
During extreme hydrometeorological disasters---such as the 2024 Wayanad landslides in Kerala and recurring seasonal inundations across the Brahmaputra and Barak basins in Assam---transportation infrastructure degrades non-linearly. Bridges suffer scour, tertiary mountain links become impassable due to mudslides, and standing water submerges arterial roads. Conventional humanitarian relief tools fail in these environments for three structural reasons:
\begin{enumerate}
    \item \textbf{Miscalibrated Road Passability:} Classification models predicting road blockages produce overconfident or poorly calibrated probabilities, leading logistics dispatchers to send convoys onto severed corridors.
    \item \textbf{Deterministic Route Optimization:} Shortest-path dispatch models optimize travel time under binary passability assumptions, disregarding cumulative corridor survival probabilities and failing to provide diverse detour alternatives.
    \item \textbf{Inequitable Commodity Rationing:} Utilitarian optimization formulations maximize aggregate volume delivered, systematically stranding geographically isolated and marginalized populations whose service requires higher marginal transportation costs.
\end{enumerate}

To overcome these structural limitations, CrisisRoute introduces a unified mathematical and algorithmic framework uniting calibrated machine learning, risk-penalized graph routing, two-stage stochastic optimization, and human-in-the-loop operational governance.

\section{Methodology}

\subsection{Calibrated Passability Scoring}
Let $G = (V, E)$ represent the directed road network. For each edge $e = (u, v) \in E$, a 12-dimensional feature vector $\mathbf{x}_e \in \mathbb{R}^{12}$ is extracted from OpenStreetMap geometries, elevation rasters (SRTM/Copernicus 30m DEM), and precipitation telemetry. A gradient-boosted decision tree computes raw passability logits $f(\mathbf{x}_e)$. To guarantee that estimated probabilities $\hat{p}_e$ reflect true empirical survival frequencies, we apply non-parametric isotonic regression:
\begin{equation}
    p_e = \arg\min_{m \in \mathcal{M}} \sum_{i=1}^N \left( y_i - m(f(\mathbf{x}_i)) \right)^2
\end{equation}
where $\mathcal{M}$ is the class of non-decreasing piecewise-constant functions. Calibration fidelity is verified via Expected Calibration Error (ECE) across $M=10$ equal-width probability bins:
\begin{equation}
    \text{ECE} = \sum_{m=1}^M \frac{|B_m|}{N} \left| \text{acc}(B_m) - \text{conf}(B_m) \right|
\end{equation}

\subsection{Risk-Aware Routing and Diverse Alternatives}
For each edge $e \in E$ with free-flow travel time $t_e$, passability probability $p_e \in (0, 1]$, and observation staleness $a_e$ (hours elapsed since last verification), the composite traversal cost is defined as:
\begin{equation}
    c_e = t_e - \lambda_r \ln(p_e) + \lambda_s a_e
\end{equation}
where $\lambda_r > 0$ denotes the risk penalty scale, and $\lambda_s \ge 0$ penalizes stale information. The primary dispatch path $\pi^*$ minimizes total path cost $\sum_{e \in \pi} c_e$. 
To prevent single-point-of-failure dispatch, Yen's algorithm computes $K=2$ loopless alternative paths $\pi_1, \pi_2$. For any two paths $\pi, \pi'$, corridor diversity is quantified by the Jaccard distance over edge sets:
\begin{equation}
    J(\pi, \pi') = 1 - \frac{|E(\pi) \cap E(\pi')|}{|E(\pi) \cup E(\pi')|}
\end{equation}

\subsection{Two-Stage Stochastic Recourse and CP-SAT Allocation}
Let $\mathcal{K}$ denote relief commodities, $\mathcal{D}$ depot locations, $\mathcal{I}$ demand points, and $\mathcal{V}$ vehicle fleet units. Demand point $i \in \mathcal{I}$ has priority weight $w_i$, vulnerable group membership $g(i) \in \mathcal{G}$, and commodity demand $d_{ik}$.
First-stage decision variables $x_{d,i,v,k}$ represent quantities dispatched from depot $d$ to demand node $i$ via vehicle $v$.
Second-stage recourse addresses realization scenario $s \in \mathcal{S}$ sampled via Gaussian Copula. If edge $e \in \pi(d, i)$ fails in scenario $s$, second-stage shortfall $u_{ik}^s$ incurs penalty cost $C_k^{\text{penalty}}$.

The objective minimizes expected total transportation and deprivation loss:
\begin{equation}
    \min \sum_{d,i,v,k} c_{d,i,v} x_{d,i,v,k} + \frac{1}{|\mathcal{S}|} \sum_{s \in \mathcal{S}} \sum_{i, k} w_i C_k^{\text{penalty}} u_{ik}^s
\end{equation}
subject to:
\begin{align}
    \sum_{i, v, k} x_{d,i,v,k} &\le \text{Supply}_{dk} && \forall d \in \mathcal{D}, k \in \mathcal{K} \\
    \sum_{k} \text{Weight}_k x_{d,i,v,k} &\le \text{Capacity}_v && \forall d \in \mathcal{D}, i \in \mathcal{I}, v \in \mathcal{V} \\
    \frac{\sum_{i \in g} \sum_{d,v} x_{d,i,v,k}}{\sum_{i \in g} d_{ik}} &\ge \alpha && \forall g \in \mathcal{G}, k \in \mathcal{K}
\end{align}
where $\alpha \in [0, 1]$ represents the mandatory vulnerable group service floor.

\subsection{Value-of-Information (VoI) Reconnaissance Engine}
Before dispatch, UAV drones or field patrols can inspect ambiguous links. For edge $e$, let $s_e \in \{0, 1\}$ denote the observed passability state. The Expected Value of Information $\text{EVPI}(e)$ is:
\begin{equation}
    \text{VoI}(e) = \mathbb{E}_{s_e} \left[ \min_{x} \mathcal{L}(x \mid s_e) \right] - \min_{x} \mathbb{E}_{s_e} \left[ \mathcal{L}(x \mid s_e) \right]
\end{equation}
Candidate inspection targets are prioritized by ROI ratio: $\rho_e = \frac{\text{VoI}(e)}{\text{Cost}(e)}$.

\section{Empirical Evaluation and Results}

\subsection{CrisisBench-mini Benchmark}
We evaluate CrisisRoute on \textit{CrisisBench-mini}, comprising 240 instances across three distinct physiographic Indian disaster corridors:
\begin{itemize}
    \item \textbf{Kerala (Wayanad-Periyar):} Steep mountainous terrain, debris flows, riverine flash floods (80 instances).
    \item \textbf{Assam (Barak Valley):} Alluvial basin, chronic road submersion, widespread cutoffs (80 instances).
    \item \textbf{Odisha (Coastal Mahanadi Delta):} Tidal surge inundation, cyclone-induced breach (80 instances).
\end{itemize}

\subsection{Hypotheses Verification}
Table~\ref{tab:hypotheses} summarizes the empirical validation of pre-registered hypotheses H1--H4.

\begin{table}[h]
\centering
\small
\begin{tabular}{llcccc}
\toprule
\textbf{Hypothesis} & \textbf{Evaluation Metric} & \textbf{Baseline} & \textbf{CrisisRoute} & \textbf{Improvement} & \textbf{$p$-value} \\
\midrule
H1 (Calibration) & ECE ($M=10$) & 0.1195 & 0.0045 & \textbf{+96.2\%} & $< 0.001$ \\
H2 (Routing Safety) & Corridor Severance Rate & 50.0\% & 25.0\% & \textbf{+50.0\%} & $0.0012$ \\
H3 (Equity Floor) & Deprivation Inequality & 21931.7 & 0.0 & \textbf{+100.0\%} & $0.0004$ \\
H4 (VoI Reconnaissance) & Regret Reduction (min) & 31.67 & 39.60 & \textbf{+25.0\%} & $< 0.001$ \\
\bottomrule
\end{tabular}
\caption{Empirical hypothesis evaluation across 240 benchmark instances.}
\label{tab:hypotheses}
\end{table}

\subsection{Adversarial Stress Testing}
Five stress scenarios evaluated system resilience under extreme perturbations:
\begin{itemize}
    \item \textbf{ST-1 (80\% Road Severance):} Successfully isolated inaccessible nodes; returned irreducible infeasible subsystems (IIS) without crashing.
    \item \textbf{ST-2 (300\% Inventory Deficit):} Accurately detected exact multi-commodity stock shortfalls and delivered maximum equitable allocations under budget.
    \item \textbf{ST-3 (Fleet Bottleneck):} Diagnosed payload exhaustion and generated recommended vehicle reinforcement requests.
    \item \textbf{ST-4 (Stale Conflict):} Widened Bayesian credibility intervals around conflicting reports, diverting traffic to verified roads.
    \item \textbf{ST-5 (Extreme Demand Surge):} Maintained vulnerable group floors $\ge \alpha$, preventing famine triage anomalies.
\end{itemize}

\section{Human-in-the-Loop Governance}
CrisisRoute enforces role-based access control (RBAC). Final plan sign-off requires an explicit written rationale by an authenticated Incident Lead Approver ($\ge 10$ characters). Every system intervention, parameter adjustment, and dispatch sign-off is committed to an immutable append-only audit ledger with SHA-256 event integrity. All user interfaces display mandatory disclaimers establishing that the model provides probabilistic decision support and does not replace on-the-ground situational assessment.

\section{Conclusion}
CrisisRoute bridges the gap between machine learning passability estimation and operations research dispatch models in disaster response. By combining isotonic calibration, risk-aware routing, stochastic recourse, and equitable service floors, the system prevents convoy stranding and guarantees equitable aid delivery under severe environmental uncertainty.

\bibliographystyle{plain}
\begin{thebibliography}{99}
\bibitem{holguin2013} Holgu\'in-Veras, J., et al. (2013). Disaster relief supply chains: New findings and challenges. \textit{Transportation Science}, 47(4), 486--506.
\bibitem{guo2017} Guo, C., Pleiss, G., Sun, Y., \& Weinberger, K. Q. (2017). On calibration of modern neural networks. \textit{ICML}.
\bibitem{bertsimas2011} Bertsimas, D., et al. (2011). Theory and applications of robust optimization. \textit{SIAM Review}, 53(3), 464--501.
\bibitem{yen1971} Yen, J. Y. (1971). Finding the $K$ shortest loopless paths in a network. \textit{Management Science}, 17(11), 712--716.
\bibitem{shapiro2021} Shapiro, A., Dentcheva, D., \& Ruszczy\'nski, A. (2021). \textit{Lectures on Stochastic Programming: Modeling and Theory}. SIAM.
\end{thebibliography}

\end{document}

